\section{Preliminaries for physical renormalization}

We first collect the transfer-map, channel, and physical-closure notions used
to state the mixed-state renormalization condition. The renormalization fixed
point of \cite[Definition~4.1]{Cirac2016MPDO_arXiv} is stated in the next
section.


% The concrete dependent-sum equivalences for T_0, S_0, T_2, and S_2 are
% stated in ch21_mpdo_rfp_simple_local_normalized_preparations.tex after the eta preparations.


\begin{theorem}[Tensor extension preserves trace-preserving complete positivity]
    \label{thm:mpdo_tensor_map_id_cptp}
    \lean{Matrix.tensorMapIdLM_isKrausCPTP}
    \lean{Matrix.idTensorMapLM_isKrausCPTP}
    \leanok

    Let $\mathcal S:\End_{\C}(H)\to
        \End_{\C}(K)$ be trace-preserving and completely positive,
    and let $R$ be another finite-dimensional space. Then
    \begin{align}
        \mathcal S\otimes\id_R:
        \End_{\C}(H\otimes R) & \longrightarrow
        \End_{\C}(K\otimes R),
        \notag                                  \\
        \id_R\otimes\mathcal S:
        \End_{\C}(R\otimes H) & \longrightarrow
        \End_{\C}(R\otimes K)
        \notag
    \end{align}
    are trace-preserving and completely positive.
\end{theorem}
\begin{proof}
    \leanok
    The first assertion is the tensor-extension closure theorem. For the
    second, let $\tau_H:R\otimes H\simeq H\otimes R$ and
    $\tau_K:K\otimes R\simeq R\otimes K$ be the canonical factor swaps. Then
    \begin{align}
        \id_R\otimes\mathcal S
         & =\mathcal R_{\tau_K}\circ
        (\mathcal S\otimes\id_R)\circ\mathcal R_{\tau_H},
        \notag
    \end{align}
    where $\mathcal R_\tau$ denotes reindexing both matrix coordinates along
    $\tau$. Both reindexing maps are trace-preserving and completely positive,
    so closure under composition proves the claim. No choice of Kraus family
    is needed for this preservation statement.
\end{proof}


\begin{lemma}[Compressed sums of completely positive maps]
    \label{lem:mpdo_compressed_cp_sum}
    \lean{isKrausCP_sum_comp_singleKrausMap}
    \leanok

    Let $(P_i)_i$ be finitely many operators on $H$ and let each
    $\Phi_i:\End_{\C}(H)\to\End_{\C}(K)$ be completely positive. Then
    \begin{align}
        X & \mapsto\sum_i\Phi_i(P_iXP_i^\dagger)
        \notag
    \end{align}
    is completely positive. No completeness or orthogonality condition on
    the $P_i$ is required.
\end{lemma}
\begin{proof}\leanok

    Composing each Kraus family of $\Phi_i$ with the corresponding $P_i$
    gives one combined rectangular Kraus family representing the sum.
\end{proof}

\begin{lemma}[Trace selection by projector-controlled sums]
    \label{lem:mpdo_projector_controlled_trace_selection}
    \lean{trace_sum_comp_singleKrausMap_of_starProjections}
    \leanok

    Let $(P_i)_i$ be finitely many orthogonal projections on $H$, not
    necessarily resolving the identity, and let each
    $\Phi_i:\End_{\C}(H)\to\End_{\C}(K)$ be trace-preserving and
    completely positive. Then, for every $X$,
    \begin{align}
        \tr\Bigl(\sum_i\Phi_i(P_iXP_i)\Bigr)
         & =\tr\Bigl(\Bigl(\sum_iP_i\Bigr)X\Bigr).
        \notag
    \end{align}
\end{lemma}
\begin{proof}\leanok
    Trace preservation of each $\Phi_i$ reduces the left side to
    $\sum_i\tr(P_iXP_i)$; cyclicity of the trace and idempotence of each
    projection give $\tr(P_iXP_i)=\tr(P_iX)$.
\end{proof}

% Classification: conditional helper.
% This is a project-derived abstraction of the projection-first channel
% assembly in CPSV16, Appendix C.2, lines 1810--1817, not a theorem stated
% there in this general form.


The state-preparation map $\mathcal P_\rho(X)=X\otimes\rho$ from matrices on
$A$ to matrices on $A\otimes B$, the elementary preparation operation used
in the maps $\mathcal T_1$ and $\mathcal S_1$ of
\cite[Appendix~C.2, lines~1527--1533 and~1551--1555]{Cirac2016MPDO_arXiv},
is a fully generic construction with no MPDO-specific content, stated as
Definition~\cite[``State-preparation map'']{QICLean2026Blueprint} in the
companion quantum-channel volume.


\begin{definition}[Physical closure maps]
    \label{def:phys_close}
    \lean{MPOTensor.physCloseN, MPOTensor.physCloseN_apply,
        MPOTensor.physClose1,
        MPOTensor.physClose2, MPOTensor.physClose3}
    \leanok
    \uses{def:mpo_eval_word}
    For an MPO tensor $M$, a length $N$, and a virtual operator $X$, define the
    physical operator $M_N(X)$ by
    \begin{align}
        M_N(X)_{\boldsymbol i,\boldsymbol j}
         & =
        \tr\!\left(M^{i_0j_0}M^{i_1j_1}\cdots M^{i_{N-1}j_{N-1}}X\right).
        \notag
    \end{align}
    This is linear in $X$. The displayed order chooses the cut of the cyclic
    virtual contraction immediately before the first site, so that $X$ follows
    the last site. The one- and two-site operators are constructed in
    \cite[lines~638--654 and Definition~4.1]{Cirac2016MPDO_arXiv}. When
    $M=\mathcal K$, write these operators as $\mathcal K_N(X)$; the cases
    $N=2,3$ occur in
    \cite[Proposition~C.7, lines~1510--1516]{Cirac2016MPDO_arXiv}.
\end{definition}

\begin{lemma}[Periodic closure]
    \label{lem:phys_close_identity}
    \lean{MPOTensor.physCloseN_identity_eq_mpo}
    \leanok
    \uses{def:phys_close, def:mpo_family}
    Closing the virtual legs with the identity gives the periodic MPO
    operator: $M_N(1)=\rho^{(N)}(M)$.
\end{lemma}

\begin{proof}
    \leanok
    For every pair of physical words $\sigma,\tau$,
    \begin{align}
        M_N(1)_{\sigma,\tau}
         & =\tr\!\left(M^{\sigma,\tau}\cdot1\right)
        =\tr\!\left(M^{\sigma,\tau}\right)
        =\rho^{(N)}(M)_{\sigma,\tau}.
        \notag
    \end{align}
\end{proof}

\begin{theorem}[One- and two-site physical closures]
    \label{thm:phys_close_one_two}
    \lean{MPOTensor.physCloseN_one_eq_physClose1,
        MPOTensor.physCloseN_two_eq_physClose2}
    \leanok
    \uses{def:phys_close}
    Under the canonical identifications of one-site configurations with
    physical indices and two-site configurations with pairs of physical
    indices,
    \begin{align}
        M_1(X)_{i,j}
         & =\tr(M^{ij}X),
        \notag                          \\
        M_2(X)_{(i_0,i_1),(j_0,j_1)}
         & =\tr(M^{i_0j_0}M^{i_1j_1}X).
        \notag
    \end{align}
\end{theorem}
\begin{proof}
    \leanok
    Under these identifications, the length-one and length-two word
    evaluations are $M^{ij}$ and $M^{i_0j_0}M^{i_1j_1}$, respectively.
\end{proof}

\begin{lemma}[Three-site physical closure]
    \label{lem:phys_close_three}
    \lean{MPOTensor.physClose3_apply, MPOTensor.physCloseN_three_eq_physClose3}
    \leanok
    \uses{def:phys_close, def:three_coordinate_identification}
    The general three-site physical closure has coefficients
    \begin{align}
        M_3(X)_{(i_0,(i_1,i_2)),(j_0,(j_1,j_2))}
         & =\tr(M^{i_0j_0}M^{i_1j_1}M^{i_2j_2}X).
        \notag
    \end{align}
    For $M=\mathcal K$, this is the operator $\mathcal K_3(X)$ in
    \cite[Proposition~C.7, lines~1510--1516]{Cirac2016MPDO_arXiv}.
\end{lemma}
\begin{proof}
    \leanok
    The length-three word evaluation is
    $M^{i_0j_0}M^{i_1j_1}M^{i_2j_2}$.
\end{proof}

\section{Renormalization fixed points}

The following is the mixed-state fixed-point condition of
\cite[Definition~4.1]{Cirac2016MPDO_arXiv}. Its two local channels extend to
every longer periodic chain by acting on the first one or two sites and leaving
the remaining sites unchanged.

\begin{definition}[Renormalization fixed point via trace-preserving maps]
    \label{def:is_rfp_via_ts}
    \lean{MPOTensor.IsRFPViaTS}
    \leanok
    \uses{def:mpo_tensor, def:phys_close}
    An MPO tensor $M$ is a \emph{renormalization fixed point} if there exist two
    trace-preserving completely positive maps $\mathcal{S}, \mathcal{T}$ on the
    physical indices such that $\mathcal{S}[M_2(X)]=M_1(X)$ and
    $\mathcal{T}[M_1(X)]=M_2(X)$ for every virtual operator $X$, where
    $M_1(X)_{ij}=\tr(M^{ij}X)$ and
    $M_2(X)_{(i_1i_2)(j_1j_2)}
        =\tr(M^{i_1j_1}M^{i_2j_2}X)$ are the one-
    and two-site physical operators obtained by closing one or two tensors with
    $X$. This condition is distinct from idempotence of the doubled-index
    transfer map for general mixed states. Its source zero-correlation-length
    consequence concerns the physical-trace transfer and is proved below.
    In \cite[Definition~4.1, lines~638--660]{Cirac2016MPDO_arXiv}, this local
    condition is imposed on a tensor already in canonical form and generating
    matrix product density operators. The definition above isolates the two
    channel equations; results that use the source's standing hypotheses state
    them separately.
    On the open coefficient tensors the two maps have types:
    % Formula: for every pair of virtual indices \alpha,\beta, set
    % (M_{\alpha\beta})_{ij}=M^{ij}_{\alpha\beta} and
    % ((MM)_{\alpha\beta})_{(i_1,i_2),(j_1,j_2)}
    % =\sum_\gamma M^{i_1j_1}_{\alpha\gamma}
    % M^{i_2j_2}_{\gamma\beta}.  The pictured identities are
    % \mathcal T(M_{\alpha\beta})=(MM)_{\alpha\beta} and
    % \mathcal S((MM)_{\alpha\beta})=M_{\alpha\beta} for every \alpha,\beta.
    % Here i,j,\alpha,\beta are free on M_{\alpha\beta}, giving signature
    % (1,1,1,1).  On (MM)_{\alpha\beta}, i_1,i_2,j_1,j_2,\alpha,\beta are
    % free and \gamma is contracted between the two copies of M, giving
    % signature (1,1,2,2).
    % Contracting \alpha,\beta against X gives
    % M_1(X)_{ij}=\tr(M^{ij}X) and
    % M_2(X)_{(i_1,i_2),(j_1,j_2)}
    % =\tr(M^{i_1j_1}M^{i_2j_2}X), hence
    % \mathcal T[M_1(X)]=M_2(X) and \mathcal S[M_2(X)]=M_1(X) for every X.
    % Their X-closed signatures are (0,0,1,1) and (0,0,2,2), respectively.
    % Conversely, nondegeneracy of the trace pairing recovers the pictured
    % open coefficient identities from the X-closed identities.
    % Ink: each square M carries one west and one east virtual index and one
    % north and one south physical index; the bond in each two-site object is
    % the contracted virtual index \gamma.  Each blue stroke is a typed map
    % between the complete objects at its endpoints, never a tensor contraction.
    \par\noindent\hfil
    {\tenkzkernel
        \begin{tikzcd}[channel, ampersand replacement=\&,
                column sep=10mm, row sep=8mm]
            {\begin{tenkz}[cols=1, west=open, east=open, physical=updown]
                    \tn[skin=mpo]{M}
                \end{tenkz}} \arrow[r, "\mathcal T"] \&
            {\begin{tenkz}[cols=2, west=open, east=open, physical=updown]
                        \tn[skin=mpo]{M} & \tn[skin=mpo]{M}
                    \end{tenkz}} \\
            {\begin{tenkz}[cols=2, west=open, east=open, physical=updown]
                \tn[skin=mpo]{M} & \tn[skin=mpo]{M}
            \end{tenkz}} \arrow[r, "\mathcal S"] \&
            {\begin{tenkz}[cols=1, west=open, east=open, physical=updown]
                        \tn[skin=mpo]{M}
                    \end{tenkz}}
        \end{tikzcd}}
    \hfil\par
\end{definition}

\begin{theorem}[Pure-state specialization of the physical-map fixed point]
    \label{thm:pure_doubled_tensor_rfp_via_ts}
    \lean{MPSTensor.doubledTensor_isRFPViaTS_iff_isTransferIdempotent,
        MPSTensor.doubledTensor_isRFPViaTS_iff_hasPhysicalBlockingIsometry}
    \leanok
    \uses{def:is_rfp_via_ts, def:mps_transfer_idempotence, def:physical_blocking_isometry}
    Let $A=(A^i)_{i=0}^{d-1}$ be an MPS tensor, and let
    $M(A)^{ij}=A^i\otimes\overline{A^j}$ be its doubled MPO tensor. Then
    \begin{align}
        M(A)\text{ satisfies Definition~\ref{def:is_rfp_via_ts}}
         & \quad\Longleftrightarrow\quad \E_A^2=\E_A
        \notag                                       \\
         & \quad\Longleftrightarrow\quad
        A\text{ has a physical blocking isometry}.
        \notag
    \end{align}
    More explicitly, if $V\colon\C^d\to\C^{d^2}$ is the blocking isometry,
    the refinement channel has the orientation
    \begin{align}
        \mathcal T(X) & =VXV^\dagger.
        \notag
    \end{align}
    If $d>0$, put $P=VV^\dagger$ and choose a density matrix $\rho_0$ on the
    one-site space. A trace-preserving coarse-graining channel is
    \begin{align}
        \mathcal S(Y)
         & =V^\dagger YV
        +\tr\!\bigl((\Id-P)Y(\Id-P)\bigr)\rho_0.
        \notag
    \end{align}
    On the range of $V$, the second term vanishes and this agrees with the
    source construction
    $\mathcal S(Y)=\tr_2(U^\dagger YU)$ after extending $V$ to a unitary with
    a fixed pure ancilla. If $d=0$, both matrix spaces are trivial and the
    unique zero maps give $\mathcal S$ and $\mathcal T$. This is the
    pure-state specialization stated at
    \cite[line~668, labels \texttt{Def-equiv-pure} and
        \texttt{T-and-S-mixed-equiv-pure}, lines~675--687]
    {Cirac2016MPDO_arXiv}.
\end{theorem}
\begin{proof}
    \leanok
    \uses{thm:physical_blocking_isometry_iff_transfer_idempotence, thm:rfp_via_ts_physical_trace_idempotent, thm:matrix_support_completion_cptp}
    The blocking relation
    $A^{i_1}A^{i_2}=\sum_jV_{(i_1,i_2),j}A^j$ implies
    \begin{align}
        M_2(X) & =V M_1(X)V^\dagger.
        \notag
    \end{align}
    For $d>0$, the refinement identity follows directly:
    \begin{align}
        \mathcal T(M_1(X))
         & =V M_1(X)V^\dagger
        =M_2(X).
        \notag
    \end{align}
    For the coarse-graining identity, put $P=VV^\dagger$. Since
    $V^\dagger V=\Id$ and $(\Id-P)V=0$, substitution gives
    \begin{align}
        \mathcal S(M_2(X))
         & =V^\dagger(VM_1(X)V^\dagger)V
        +\tr\!\bigl((\Id-P)(VM_1(X)V^\dagger)(\Id-P)\bigr)\rho_0
        \notag                           \\
         & =M_1(X)+0
        =M_1(X).
        \notag
    \end{align}
    If $d=0$, the unique zero channels satisfy both identities. Conversely,
    trace preservation of $\mathcal T$ gives idempotence of the physical-trace
    transfer of $M(A)$, as in
    \cite[Appendix C, lines~1333--1340]{Cirac2016MPDO_arXiv}. After exchanging
    the two tensor factors, this matrix is the transfer matrix of $\E_A$.
    Faithfulness of the transfer-matrix representation therefore gives
    $\E_A^2=\E_A$. The physical-isometry formulation follows from
    Theorem~\ref{thm:physical_blocking_isometry_iff_transfer_idempotence}.
\end{proof}
% CPSV16 Appendix D, lines 2116--2176.  The labels `rank-Fibonacci` and
% `eq:1` refer respectively to the rank formula and its internal calculation.
\begin{definition}[Fibonacci vacuum-sector support tensor]
    \label{def:cpsv_fibonacci_support_tensor}
    \lean{MPSTensor.FibonacciBoundary.ExactlyTwoZero,
        MPSTensor.FibonacciBoundary.FusionWeights,
        MPSTensor.FibonacciBoundary.canonicalFusionWeights,
        MPSTensor.FibonacciBoundary.physicalTriple,
        MPSTensor.FibonacciBoundary.tensor,
        MPSTensor.FibonacciBoundary.canonicalTensor}
    \leanok
    \uses{def:mps_tensor}
    This is the concrete tensor in
    \cite[Appendix~D, lines~2116--2121]{Cirac2016MPDO_arXiv}.  Its physical
    labels are triples $i,j,k\in\{0,1\}$ and its virtual labels are
    $\alpha,\beta\in\{0,1\}$.  The fusion weights satisfy
    \begin{align}
        N_{ijk}=0
         & \iff (i,j,k)\in\{(0,0,1),(0,1,0),(1,0,0)\}.
        \label{eq:cpsv_fib_support}
    \end{align}
    Every weight outside this zero set is positive, and
    \begin{align}
        A^{ijk}_{\alpha\beta}
         & =\delta_{i,\alpha}\delta_{k,\beta}N_{ijk}.
        \notag
    \end{align}
    The canonical support tensor takes every nonzero $N_{ijk}$ to be one.
\end{definition}

\begin{definition}[Fibonacci transition matrix]
    \label{def:cpsv_fibonacci_transition_matrix}
    \lean{MPSTensor.FibonacciBoundary.B}
    \leanok
    The transition matrix from
    \cite[Appendix~D, lines~2146--2149]{Cirac2016MPDO_arXiv} is
    \begin{align}
        B & =\begin{pmatrix}1&1\\1&2\end{pmatrix}.
        \notag
    \end{align}
\end{definition}

\begin{definition}[Fibonacci boundary transition counts]
    \label{def:cpsv_fibonacci_boundary_counts}
    \lean{MPSTensor.FibonacciBoundary.x}
    \leanok
    \uses{def:cpsv_fibonacci_transition_matrix}
    For $n\in\N$ and $\alpha,\beta\in\{0,1\}$, define
    \begin{align}
        x^n_{\alpha\beta} & =(B^n)_{\alpha\beta}.
        \label{eq:cpsv_fib_counts}
    \end{align}
    These are the transition counts governed by the recurrence in
    \cite[Appendix~D, lines~2131--2149]{Cirac2016MPDO_arXiv}.
\end{definition}

\begin{theorem}[Fibonacci boundary transition recurrence]
    \label{thm:cpsv_fibonacci_boundary_recurrence}
    \lean{MPSTensor.FibonacciBoundary.x_one,
        MPSTensor.FibonacciBoundary.x_succ}
    \leanok
    \uses{def:cpsv_fibonacci_boundary_counts}
    In the order $(00,01,10,11)$, one has $x^1=(1,1,1,2)$ and
    \begin{align}
        x^{n+1}_{\alpha0} & =x^n_{\alpha0}+x^n_{\alpha1},  &
        x^{n+1}_{\alpha1} & =x^n_{\alpha0}+2x^n_{\alpha1}.
        \notag
    \end{align}
    This is the recurrence in
    \cite[Appendix~D, lines~2131--2145]{Cirac2016MPDO_arXiv}.
\end{theorem}

\begin{proof}
    \leanok
    \uses{def:cpsv_fibonacci_transition_matrix, def:cpsv_fibonacci_boundary_counts}
    Multiply the row indexed by $\alpha$ in $B^n$ by the two columns of $B$.
\end{proof}

\begin{definition}[Periodic Fibonacci transition count]
    \label{def:cpsv_fibonacci_periodic_count}
    \lean{MPSTensor.FibonacciBoundary.periodicTransitionCount}
    \leanok
    \uses{def:cpsv_fibonacci_boundary_counts}
    Closing the two boundary labels gives the transition count
    \begin{align}
        a_N & =x^N_{00}+x^N_{11}=\tr(B^N).
        \notag
    \end{align}
    This is the boundary-label closure used in
    \cite[Appendix~D, line~2157]{Cirac2016MPDO_arXiv}.
\end{definition}

\begin{theorem}[Characteristic recurrence for periodic Fibonacci transition counts]
    \label{thm:cpsv_fibonacci_periodic_count_recurrence}
    \lean{MPSTensor.FibonacciBoundary.periodicTransitionCount_add_two_add}
    \leanok
    \uses{def:cpsv_fibonacci_periodic_count}
    For every $N\in\N$,
    \begin{align}
        a_{N+2}+a_N & =3a_{N+1}.
        \notag
    \end{align}
\end{theorem}

\begin{proof}
    \leanok
    \uses{def:cpsv_fibonacci_transition_matrix, def:cpsv_fibonacci_boundary_counts, def:cpsv_fibonacci_periodic_count}
    Expand $B^{N+2}=B^{N+1}B$ and $B^{N+1}=B^NB$. Substitution of
    $B=\left(\begin{smallmatrix}1&1\\1&2\end{smallmatrix}\right)$ into the
    diagonal sums gives $a_{N+2}+a_N=3a_{N+1}$.
\end{proof}

\begin{theorem}[First periodic Fibonacci transition counts]
    \label{thm:cpsv_fibonacci_periodic_count_values}
    \lean{MPSTensor.FibonacciBoundary.periodicTransitionCount_one,
        MPSTensor.FibonacciBoundary.periodicTransitionCount_two}
    \leanok
    \uses{def:cpsv_fibonacci_periodic_count, thm:cpsv_fibonacci_boundary_recurrence}
    The first two periodic transition counts are $a_1=3$ and $a_2=7$.
\end{theorem}

\begin{proof}
    \leanok
    \uses{thm:cpsv_fibonacci_boundary_recurrence}
    Substitute $x^1=(1,1,1,2)$ into the recurrence once and close the two
    diagonal boundary labels.
\end{proof}

\begin{theorem}[Golden-ratio formula for periodic Fibonacci transition counts]
    \label{thm:cpsv_fibonacci_periodic_count_golden_ratio}
    \lean{MPSTensor.FibonacciBoundary.periodicTransitionCount_eq_goldenRatio}
    \leanok
    \uses{def:cpsv_fibonacci_periodic_count}
    Set $\tau_+=(1+\sqrt5)/2$ and $\tau_-=(1-\sqrt5)/2$. For every
    $N\in\N$,
    \begin{align}
        a_N & =\tau_+^{2N}+\tau_-^{2N}.
        \notag
    \end{align}
    This is the transition-count formula in
    \cite[Appendix~D, lines~2146--2176]{Cirac2016MPDO_arXiv}.
\end{theorem}

\begin{proof}
    \leanok
    \uses{def:cpsv_fibonacci_transition_matrix, def:cpsv_fibonacci_boundary_counts, def:cpsv_fibonacci_periodic_count, thm:cpsv_fibonacci_periodic_count_recurrence}
    Use two-step induction on $N$. The cases $N=0$ and $N=1$ follow by direct
    calculation. Since $\tau_\pm^2=\tau_\pm+1$, each sequence
    $\tau_\pm^{2N}$ satisfies
    $\tau_\pm^{2(N+2)}+\tau_\pm^{2N}=3\tau_\pm^{2(N+1)}$. The induction step
    now follows from the same recurrence for $a_N$.
\end{proof}

\begin{theorem}[Periodic Fibonacci transition counts are not geometric]
    \label{thm:cpsv_fibonacci_transition_not_geometric}
    \lean{MPSTensor.FibonacciBoundary.periodicTransitionCount_not_geometric}
    \leanok
    \uses{def:cpsv_fibonacci_periodic_count, thm:cpsv_fibonacci_periodic_count_values}
    There are no $r,s\in\N$ such that $a_N=rs^{N-1}$ for every $N>0$.
    This is the non-geometricity assertion at
    \cite[Appendix~D, line~2125]{Cirac2016MPDO_arXiv} for the transition-count
    sequence.
\end{theorem}

\begin{proof}
    \leanok
    \uses{thm:cpsv_fibonacci_periodic_count_values}
    Such a formula would give $r=a_1=3$ and $7=a_2=3s$, which is impossible
    in $\N$.
\end{proof}

\begin{definition}[Open-boundary Fibonacci operator]
    \label{def:cpsv_fibonacci_open_boundary_operator}
    \lean{MPSTensor.FibonacciBoundary.openBoundaryOperator,
        MPSTensor.FibonacciBoundary.openBoundaryOperator_apply}
    \leanok
    \uses{def:cpsv_fibonacci_support_tensor}
    For $n>0$ and $\alpha,\beta\in\{0,1\}$, define the matrix
    $\rho^{(n)}_{\alpha\beta}(A)$, indexed by physical words $\sigma$ and $\tau$, by
    \begin{align}
        \bigl(\rho^{(n)}_{\alpha\beta}(A)\bigr)_{\sigma,\tau}
         & =\begin{cases}
                (\alpha|A^{\sigma_1}\cdots A^{\sigma_n}|\beta), & \sigma=\tau,   \\
                0,                                              & \sigma\ne\tau.
            \end{cases}
    \end{align}
\end{definition}

\begin{theorem}[Open-boundary Fibonacci operator ranks]
    \label{thm:cpsv_fibonacci_open_boundary_rank}
    \lean{MPSTensor.FibonacciBoundary.rank_openBoundaryOperator_eq_x}
    \leanok
    \uses{def:cpsv_fibonacci_open_boundary_operator, def:cpsv_fibonacci_boundary_counts}
    For $n>0$ and $\alpha,\beta\in\{0,1\}$,
    \begin{align}
        \operatorname{rank}\rho^{(n)}_{\alpha\beta}(A)
         & =x^n_{\alpha\beta}.
    \end{align}
    This is \cite[Appendix~D, lines~2127--2130]{Cirac2016MPDO_arXiv}.
\end{theorem}

\begin{proof}
    \leanok
    \uses{def:cpsv_fibonacci_open_boundary_operator, def:cpsv_fibonacci_transition_matrix, def:cpsv_fibonacci_boundary_counts}
    A diagonal entry is nonzero exactly when its physical labels form an allowed
    path from $\alpha$ to $\beta$. Splitting a path at its first physical label
    shows that the numbers of supported words obey the transition rule determined
    by $B$, with the same length-zero initial values as $x^n_{\alpha\beta}$.
    Induction therefore identifies the number of nonzero diagonal entries with
    $x^n_{\alpha\beta}$. The rank of a diagonal matrix is the number of its nonzero
    diagonal entries.
\end{proof}

\begin{theorem}[Fibonacci periodic operator-rank formula]
    \label{thm:cpsv_fibonacci_periodic_rank}
    \lean{MPSTensor.FibonacciBoundary.rank_mpo_toMPOTensor_eq_periodicTransitionCount,
        MPSTensor.FibonacciBoundary.rank_mpo_toMPOTensor_eq_goldenRatio}
    \leanok
    \uses{def:cpsv_fibonacci_support_tensor, def:cpsv_fibonacci_periodic_count}
    Let $A$ be obtained from any positive choice of fusion weights with the
    support prescribed in Definition~\ref{def:cpsv_fibonacci_support_tensor}.
    Set $\tau_\pm=(1\pm\sqrt5)/2$, as in
    \cite[Appendix~D, line~2125]{Cirac2016MPDO_arXiv}. For every $N>0$,
    \begin{align}
        \operatorname{rank}\rho^{(N)}(A)
         & =a_N=x^N_{00}+x^N_{11}
        =\tau_+^{2N}+\tau_-^{2N}.
        \label{eq:cpsv_fib_periodic_rank}
    \end{align}
    This is \cite[Appendix~D, lines~2122--2125 and
        2157--2176]{Cirac2016MPDO_arXiv}.
\end{theorem}

\begin{proof}
    \leanok
    \uses{def:cpsv_fibonacci_open_boundary_operator, def:cpsv_fibonacci_periodic_count, thm:cpsv_fibonacci_periodic_count_golden_ratio, thm:cpsv_fibonacci_open_boundary_rank, thm:to_mpo_diagonal_rank}
    The diagonal MPO rank counts physical words with nonzero periodic amplitude.
    For $N>0$, the supports obtained by closing the virtual endpoint at $0$ and
    at $1$ are disjoint, since the first physical bit determines the initial
    endpoint. Their cardinalities are $x^N_{00}$ and $x^N_{11}$ by the
    open-boundary formula. The golden-ratio identity for $a_N$ gives the final
    equality.
\end{proof}

\begin{theorem}[Fibonacci periodic operator ranks are not geometric]
    \label{thm:cpsv_fibonacci_operator_rank_not_geometric}
    \lean{MPSTensor.FibonacciBoundary.rank_mpo_toMPOTensor_not_geometric}
    \leanok
    \uses{def:cpsv_fibonacci_support_tensor}
    Let $A$ be obtained from any positive choice of fusion weights with the
    support prescribed in Definition~\ref{def:cpsv_fibonacci_support_tensor}.
    There are no $r,s\in\N$ such that
    $\operatorname{rank}\rho^{(N)}(A)=rs^{N-1}$ for every $N>0$.
    This is the operator-level conclusion in
    \cite[Appendix~D, line~2125]{Cirac2016MPDO_arXiv}.
\end{theorem}

\begin{proof}
    \leanok
    \uses{thm:cpsv_fibonacci_periodic_rank, thm:cpsv_fibonacci_transition_not_geometric}
    The operator rank equals $a_N$ at every positive length, while
    $a_1=3$ and $a_2=7$. A geometric formula would therefore give $r=3$ and
    $7=3s$, which is impossible in $\N$.
\end{proof}

\begin{theorem}[The Fibonacci boundary tensor is not a strong fixed point]
    \label{thm:cpsv_fibonacci_not_strong_rfp}
    \lean{MPSTensor.FibonacciBoundary.toMPOTensor_not_isStrongRFP}
    \leanok
    \uses{def:cpsv_fibonacci_support_tensor, def:mpdo_strong_rfp}
    Let $A$ be obtained from any positive choice of fusion weights with the
    support prescribed in Definition~\ref{def:cpsv_fibonacci_support_tensor}.
    Then its canonical diagonal MPO tensor is not a strong renormalization
    fixed point.
\end{theorem}

\begin{proof}
    \leanok
    \uses{thm:cpsv_fibonacci_operator_rank_not_geometric, thm:mpdo_strong_rfp_periodic_rank_growth}
    Suppose otherwise, and let $P$ be the positive witness for the strong
    fixed-point relation. Then, for every $N>0$, the periodic ranks satisfy
    \begin{align}
        \operatorname{rank}\rho^{(N)}(A)
         & =\operatorname{rank}\rho^{(1)}(A)
        \bigl(\operatorname{rank}P\bigr)^{N-1}.
        \label{eq:cpsv_fib_strong_rfp_rank}
    \end{align}
    (\ref{eq:cpsv_fib_strong_rfp_rank}) contradicts the
    non-geometricity of the Fibonacci periodic ranks.
\end{proof}

\begin{definition}[Initial-coordinate regroupings]
    \label{def:finite_tuple_regroupings}
    \lean{finSuccArrowEquiv,
        finAddTwoArrowEquiv}
    \leanok
    Let $R_1$ identify an $(N+1)$-tuple with its first coordinate and the
    remaining $N$-tuple. Let $R_2$ identify an $(N+2)$-tuple with its first
    two coordinates and the remaining $N$-tuple.
\end{definition}

\begin{definition}[Renormalization maps on a longer ring]
    \label{def:mpdo_global_renormalization_maps}
    \lean{MPOTensor.refineFirstSite,
        MPOTensor.coarsenFirstTwoSites}
    \leanok
    \uses{def:finite_tuple_regroupings}
    Let $N$ count the sites that are left unchanged. For physical maps
    $\mathcal T:M_d\to M_{d^2}$ and $\mathcal S:M_{d^2}\to M_d$, define
    \begin{align}
        \widehat{\mathcal T}_N
         & =\mathcal T\otimes\id_{d^N},
        \notag                          \\
        \widehat{\mathcal S}_N
         & =\mathcal S\otimes\id_{d^N},
        \notag
    \end{align}
    using $R_1$ and $R_2$. Thus $N=0$ leaves no spectator site, but the
    maps still act on a one-site or two-site operator; no empty ring occurs.
\end{definition}

\begin{lemma}[Regrouping the first physical sites]
    \label{lem:mpdo_global_renormalization_regrouping}
    \lean{finSuccArrowEquiv_apply,
        finSuccArrowEquiv_symm_apply,
        finAddTwoArrowEquiv_apply,
        finAddTwoArrowEquiv_symm_apply}
    \leanok
    \uses{def:finite_tuple_regroupings}
    For a physical word $\sigma$ of length $N+1$,
    \begin{align}
        R_1(\sigma)
         & =(\sigma_0,(\sigma_1,\ldots,\sigma_N)),
        \notag                                     \\
        R_1^{-1}(i,(a_1,\ldots,a_N))
         & =(i,a_1,\ldots,a_N).
        \notag
    \end{align}
    For a physical word $\sigma$ of length $N+2$,
    \begin{align}
        R_2(\sigma)
         & =\bigl((\sigma_0,\sigma_1),
        (\sigma_2,\ldots,\sigma_{N+1})\bigr),
        \notag                         \\
        R_2^{-1}((i,j),(a_1,\ldots,a_N))
         & =(i,j,a_1,\ldots,a_N).
        \notag
    \end{align}
\end{lemma}

\begin{proof}
    \leanok
    All four identities hold by definition of the product-index equivalences.
\end{proof}

\begin{lemma}[Localized renormalization maps are channels]
    \label{lem:mpdo_global_renormalization_maps_cptp}
    \lean{MPOTensor.refineFirstSite_isKrausCPTP,
        MPOTensor.coarsenFirstTwoSites_isKrausCPTP}
    \leanok
    \uses{def:mpdo_global_renormalization_maps}
    If $\mathcal S$ and $\mathcal T$ are trace-preserving completely positive,
    then so are $\widehat{\mathcal S}_N$ and
    $\widehat{\mathcal T}_N$ for every $N$.
\end{lemma}

\begin{proof}
    \leanok
    If $(K_j)_j$ are Kraus operators for $\mathcal T$ with
    $\sum_j K_j^\dagger K_j=I_d$, then the localized map has Kraus operators
    $K_j\otimes I_{d^N}$, and
    \begin{align}
        \sum_j(K_j\otimes I_{d^N})^\dagger(K_j\otimes I_{d^N})
         & =\left(\sum_j K_j^\dagger K_j\right)\otimes I_{d^N}
        =I_d\otimes I_{d^N}.
        \notag
    \end{align}
    The two regroupings are unitary conjugations and preserve complete
    positivity and the trace. The argument for $\mathcal S$ is identical.
\end{proof}

\begin{theorem}[Global action of the renormalization maps]
    \label{thm:mpdo_global_renormalization_maps}
    \lean{MPOTensor.refineFirstSite_physCloseN,
        MPOTensor.coarsenFirstTwoSites_physCloseN,
        MPOTensor.refineFirstSite_mpo,
        MPOTensor.coarsenFirstTwoSites_mpo,
        MPOTensor.exists_global_renormalization_maps}
    \leanok
    \uses{def:phys_close, def:is_rfp_via_ts, def:mpdo_global_renormalization_maps}
    Suppose that $\mathcal S[M_2(X)]=M_1(X)$ and
    $\mathcal T[M_1(X)]=M_2(X)$ for every virtual operator $X$. Then, for
    every $N$ and $X$,
    \begin{align}
        \widehat{\mathcal S}_N[M_{N+2}(X)]
         & =M_{N+1}(X),
        \notag          \\
        \widehat{\mathcal T}_N[M_{N+1}(X)]
         & =M_{N+2}(X).
        \notag
    \end{align}
    In particular, these identities hold for every periodic MPO operator.
    Hence the two maps supplied by a renormalization fixed point act on rings
    of every length as required in the proof of
    \cite[Appendix~C, lines~1333--1341]{Cirac2016MPDO_arXiv}.
\end{theorem}

\begin{proof}
    \leanok
    \uses{lem:mpdo_global_renormalization_regrouping, lem:phys_close_identity}
    For words $u,v$ on the $N$ spectator sites, write $M^{u,v}$ for their
    virtual word evaluation and write $[Y]^{(r)}_{u,v}$ for the block of an
    operator $Y$ on the first $r$ sites. Regrouping the first site gives
    \begin{align}
        [M_{N+1}(X)]^{(1)}_{u,v}
         & =M_1(M^{u,v}X).
        \notag
    \end{align}
    Therefore
    \begin{align}
        \mathcal T([M_{N+1}(X)]^{(1)}_{u,v})
         & =M_2(M^{u,v}X)
        =[M_{N+2}(X)]^{(2)}_{u,v}.
        \notag
    \end{align}
    Equality of all blocks proves the refinement identity. Interchanging
    $M_1$, $M_2$ and using $\mathcal S[M_2(Y)]=M_1(Y)$ proves the coarsening
    identity. Setting $X=1$ and using
    $M_N(1)=\rho^{(N)}(M)$
    (Lemma~\ref{lem:phys_close_identity}) gives the periodic-ring identities
    \begin{align}
        \widehat{\mathcal T}_N[\rho^{(N+1)}(M)]
         & =\rho^{(N+2)}(M),
        \notag               \\
        \widehat{\mathcal S}_N[\rho^{(N+2)}(M)]
         & =\rho^{(N+1)}(M).
        \notag
    \end{align}
\end{proof}

\section{Pure-state recovery inside the MPO formalism}

\begin{definition}[Diagonal pure-state embedding as an MPO]\label{def:to_mpo_tensor}
    \lean{MPSTensor.toMPOTensor}
    \leanok
    \uses{def:mps_tensor, def:mpo_tensor}
    An MPS tensor $A=\{A^i\}_{i=0}^{d-1}$ determines an MPO tensor by
    placing $A^i$ on the diagonal in the bra and ket indices:
    \begin{align}
        M^{ij} & =\delta_{ij}A^i.
        \notag
    \end{align}
    Equivalently, $M^{ii}=A^i$ and $M^{ij}=0$ for $i\ne j$.
\end{definition}

\begin{theorem}[Operator and rank of the diagonal pure-state embedding]
    \label{thm:to_mpo_diagonal_rank}
    \lean{MPSTensor.evalWord_toMPOTensor,
        MPSTensor.mpo_toMPOTensor,
        MPSTensor.rank_mpo_toMPOTensor}
    \leanok
    \uses{def:to_mpo_tensor, def:mpo_family, def:mpv}
    For every $N$, the embedded MPO is diagonal with diagonal
    $V^{(N)}(A)$, and hence
    \begin{align}
        \rho^{(N)}(A^{\mathrm{MPO}})
         & =\operatorname{diag}V^{(N)}(A),      &
        \operatorname{rank}\rho^{(N)}(A^{\mathrm{MPO}})
         & =\#\{\sigma:V^{(N)}(A)_\sigma\ne0\}.
        \notag
    \end{align}
\end{theorem}

\begin{proof}
    \leanok
    \uses{def:to_mpo_tensor, def:mpo_family, def:mpv}
    A ket--bra word evaluation vanishes at the first unequal pair of physical
    labels.  Equal words recover the original MPS word evaluation.  The rank
    formula follows by counting the nonzero entries of a diagonal matrix.
\end{proof}

\begin{theorem}[Transfer map of the diagonal pure-state embedding]%
    \label{thm:to_mpo_transfer_map}
    \lean{MPSTensor.toMPOTensor_transferMap}
    \leanok
    \uses{def:to_mpo_tensor, def:mpo_transfer_map, def:transfer_map}
    The transfer map of the diagonal MPO associated to $A$ agrees with the
    original MPS transfer map: $\E_{A^{\mathrm{MPO}}}=\E_A$.
\end{theorem}

\begin{proof}
    \leanok
    \uses{def:to_mpo_tensor, def:mpo_transfer_map, def:transfer_map}
    In the double sum defining $\E_{A^{\mathrm{MPO}}}$, all off-diagonal terms
    vanish because $M^{ij}=0$ for $i\ne j$, while the diagonal terms are
    exactly $A^iX(A^i)^\dagger$. Thus only the sum over $i$ remains, which
    is the defining formula for $\E_A$.
\end{proof}

\begin{theorem}[Pure-state recovery of transfer-map idempotence]%
    \label{thm:to_mpo_zcl_iff_transfer_idempotent}
    \lean{MPSTensor.toMPOTensor_isZCL_iff_isTransferIdempotent}
    \leanok
    \uses{def:mpo_is_zcl, def:mps_transfer_idempotence, thm:to_mpo_transfer_map}
    For an MPS tensor $A$, the diagonal MPO embedding has idempotent transfer
    map if and only if the MPS transfer map is idempotent:
    \begin{align}
        \E_{A^{\mathrm{MPO}}}\circ\E_{A^{\mathrm{MPO}}}
        =\E_{A^{\mathrm{MPO}}}
         & \iff
        \E_A\circ\E_A=\E_A.
        \notag
    \end{align}
\end{theorem}

\begin{proof}
    \leanok
    \uses{def:mpo_is_zcl, def:mps_transfer_idempotence, thm:to_mpo_transfer_map}
    Both conditions assert idempotence of the corresponding transfer map.
    Theorem~\ref{thm:to_mpo_transfer_map} identifies those transfer maps, so
    the two predicates are equivalent.
\end{proof}

\begin{theorem}[Pure-state recovery of zero correlation length]%
    \label{thm:to_mpo_zcl_iff_zcl}
    \lean{MPSTensor.toMPOTensor_isZCL_iff_mps_isZCL}
    \leanok
    \uses{thm:to_mpo_zcl_iff_transfer_idempotent, thm:zcl_iff_idempotent_transfer}
    For an MPS tensor $A$, the diagonal MPO embedding has idempotent transfer
    map if and only if $A$ has zero correlation length.
\end{theorem}

\begin{proof}
    \leanok
    \uses{thm:to_mpo_zcl_iff_transfer_idempotent, thm:zcl_iff_idempotent_transfer}
    Combine Theorem~\ref{thm:to_mpo_zcl_iff_transfer_idempotent} with the characterization of
    zero correlation length by idempotence of the MPS transfer map
    (Theorem~\ref{thm:zcl_iff_idempotent_transfer}).
\end{proof}
